\section{Forward models}
\label{sec:forward}

The data-assimilation pipeline is solver-agnostic: the smoother
interacts with the forward model only through an ensemble interface
that takes a parameter dataset and an optional warm-start state, and
returns a forecast trajectory.  Three solvers are wired up in
\pyurbanair{} and can be selected independently for the truth and
the assimilation model from the command line of the rollout script.

\subsection{Lattice-Boltzmann solver (\pylbm{})}
\pylbm{} is a custom lattice-Boltzmann
implementation~\cite{succi2001lbm} on a regular Cartesian grid, with
an STL-defined urban geometry and velocity-/pressure-driven
inflow--outflow boundary conditions.  In the experiments reported
here we use a moderate grid of $50 \times 40 \times 16$ cells over a
$50\,\mathrm{m}\times 40\,\mathrm{m}\times 40\,\mathrm{m}$ domain
covering the staggered cube array of Xie and
Castro~\cite{xie2008castro}.  Among the three solvers it is the
cheapest to run and is therefore the default choice for ensemble
experiments and development.

\subsection{uDALES (\pyudales{})}
uDALES~\cite{suter2018udales} is a finite-difference \LES{} code
targeted at urban-microclimate applications.  In \pyurbanair{} it is
driven through a Python wrapper that prepares the case directory,
configures a uniform or power-law vertical inflow profile, and
triggers the simulation through a MATLAB-based preprocessing chain.
Inflow nudging is applied with a fixed time constant and width to
relax the interior of the domain toward the prescribed inflow
profile.

\subsection{PALM (\pypalm{})}
PALM~\cite{maronga2020palm} is a comprehensive \LES{} model system
that supports thermodynamics, building-resolved geometry through STL
meshes, and a flexible I/O system.  Its integration in \pyurbanair{}
mirrors that of \pyudales{}: a wrapper sets up the case directory,
writes inflow and nudging configuration, and dispatches the run.

\subsection{Common interface}
Each solver exposes three Python objects with consistent signatures:
\begin{enumerate}
  \item a \emph{single-member} forward model that accepts a parameter
        dataset on a discrete time grid and integrates the flow over
        a window, optionally warm-starting from a previous terminal
        state;
  \item an \emph{ensemble} forward model that runs $\Ne$ members in
        parallel processes and gathers their outputs into a single
        ensemble dataset; and
  \item an \emph{observation operator} that samples ensemble
        velocity fields at a fixed set of probe locations and
        aggregates them in time according to the observation
        configuration.
\end{enumerate}
This common interface is what allows the rollout script to swap the
truth and the assimilation solvers independently and to remain
agnostic to the parameter time-series model in use.

\paragraph{Spin-up and warm-starting.}
Window $0$ runs with spin-up enabled to wash out the artificial
initial condition.  From window $1$ onward each ensemble member
warm-starts from the terminal state of the previous window's
forecast, so the physical state across windows is continuous up to
the discretization of the parameter trajectory.

\paragraph{Failure handling.}
Individual ensemble runs occasionally fail (e.g.\ a numerical
blow-up triggered by an outlier parameter draw).  The ensemble layer
in \pyurbanair{} handles this through a configurable failure
policy: by default, a failed member is replaced by a clone of a
randomly chosen successful member with its parameters jittered by a
small fraction of the prior spread.  This keeps the ensemble size
constant across windows and prevents a single failure from aborting
the rollout.
